% ID: SEL-D0067-OB-E01
% ESTADO: cotejada. Segunda pasada visual y borrador compilado revisados.
% Fuente: Selectividad/sel_2009_matematicas_aplicadas/mat apli CC SS01.pdf
% SHA-256: f5e34036805266809f1b5a2c6087d7ae598afd929b5179ec5ea3888db6cbccf4
% Página: 2; opción B; ejercicio original 1.

En un examen de Matemáticas se propone el siguiente problema:

``Indique dónde se alcanza el mínimo de la función $F(x,y)=6x+3y-2$ en la región determinada por las restricciones $2x+y\geq6$; $2x+5y\leq30$; $2x-y\leq6$.''
\begin{enumerate}[label=\alph*)]
\item \textbf{(2.5 puntos)} Resuelva el problema.
\item \textbf{(0.5 puntos)} Ana responde que se alcanza en $(1,4)$ y Benito que lo hace en $(3,0)$. \textquestiondown{}Es cierto que el mínimo se alcanza en $(1,4)$? \textquestiondown{}Es cierto que se alcanza en $(3,0)$?
\end{enumerate}
